\section{¿Por qué la inferencia es un sistema de ingeniería completo}

Los cursos anteriores de CS25 discutieron ampliamente las capacidades del modelo, la escala de entrenamiento y las nuevas arquitecturas; esta lección desplaza el enfoque al lugar donde los usuarios invocan realmente al modelo. La pregunta central no es tan estrecha como «cómo meter un \textbf{checkpoint} en una API», sino: ¿cómo traducir los requisitos de la aplicación a una \textbf{workload} medible y a un \textbf{contracto de corrección}, e ir implementando capa por capa a través del modelo, el motor, las GPU, la red, los contenedores, el monitoreo y la optimización? Un error de definición en cualquier capa de esta cadena puede hacer que los kernels más rápidos sean trabajo inútil.

\lecturefigure{slide-002.jpg}{De entrenamiento a inferencia: el modelo finalmente debe ser usado por personas}{Diapositiva oficial 2; Clase 00:01:55--00:02:42}

\begin{importantbox}{Definición de trabajo de esta lección sobre la inferencia}
Un servicio de inferencia no es una sola propagación hacia adelante, sino un sistema en línea o por lotes que recibe continuamente solicitudes, gestiona tokens y \textbf{KV cache}, agenda el trabajo de CPU/GPU, cumple con objetivos de latencia y confiabilidad, registra pruebas depurables y controla costos. El \textbf{forward} del modelo es solo una parte del camino de datos; los SLO, las colas, la caché, la recuperación ante fallos y la evaluación también pertenecen al sistema semántico.
\end{importantbox}

\teachervoice{00:00:56--00:01:55, el docente dice que la investigación en entrenamiento responde «de dónde viene el modelo», mientras que la inferencia responde «cómo se usa finalmente el modelo». Incluso si los lectores solo desean hacer entrenamiento, deben entender cómo su arquitectura, \textbf{tokenizer}, precisión y selección de paralelismo restringirán el posterior \textbf{serving}.}

\subsection{De centro de costos a sistema entregable}

El entrenamiento convierte capital, datos y potencia computacional en pesos, pero los pesos por sí mismos no forman automáticamente un producto cobrizable, confiable y mantenible. La inferencia envuelve los pesos en interfaces de interacción, flujos de negocio y mecanismos de confiabilidad; al mismo tiempo, el \textbf{post-training} moderno llamará extensivamente a las políticas actuales para generar despliegues (\textbf{rollout}) y devolverá los resultados al proceso de optimización, por lo que «entrenamiento» e «inferencia» ya no son dos infraestructuras completamente separadas.

\lecturefigure{slide-003.jpg}{¿Por qué la inferencia atraerá más recursos de ingeniería}{Diapositiva oficial 3; Clase 00:02:42--00:05:59}

\begin{knowledgebox}{Lectura de la figura: cuatro razones que en realidad son una cadena causal}
Primero, veamos el nivel económico: el entrenamiento consume principalmente presupuesto, mientras que el servicio asume el valor del usuario; segundo, a nivel organizacional, pocas instituciones pueden realizar preentrenamiento a gran escala, pero casi todos los productos de IA deben ejecutar inferencia; tercero, esta conecta el \textbf{post-training} de vuelta con la inferencia; y finalmente, esto indica que la inferencia abarca aplicaciones, programadores (\textbf{scheduler}), álgebra, red y hardware. La conclusión pedagógica que surge no es «el entrenamiento no es importante», sino: las capacidades del modelo solo se convierten en una experiencia estable dentro de un sistema transversal a capas.
\end{knowledgebox}

\begin{warningbox}{No se debe entender el \textbf{revenue center} como «la inferencia siempre genera ganancias»}
La inferencia solo puede crear una economía unitaria positiva cuando la calidad, la utilización, el precio y la confiabilidad se cumplen simultáneamente. GPUs con baja utilización, sobreconfiguración, reintentos de fallos, salidas largas y bucles de agentes no controlados consumirán rápidamente los ingresos; aquí se discute la ubicación de entrega de valor, no una garantía financiera.
\end{warningbox}

\subsection{Orígenes de la experiencia y límites de la evidencia}

La experiencia en producción es muy valiosa, pero las recomendaciones empíricas deben señalar su alcance de aplicación. La perspectiva del conferenciantes proviene de la plataforma \textbf{serverless GPU} de Modal, soporte a despliegues de clientes en cientos o miles de GPUs, trabajo interno de \textbf{profiling}/métricas y la propia ejecución de servicios de inferencia. Esto explica por qué esta lección enfatiza especialmente el tráfico estallido (\textbf{bursty traffic}), el inicio en frío (\textbf{cold start}), la sobrecarga del anfitrión (\textbf{host overhead}) y la \textbf{observabilidad}, e advierte a los lectores para no tomar las estadísticas de una sola plataforma como leyes físicas de todos los centros de datos.

\lecturefigure{slide-004.jpg}{Contexto de experiencia en producción de inferencia del conferenciantes}{Diapositiva oficial 4; Clase 00:06:00--00:07:35}

\begin{warningbox}{Límites de la fuente}
Este material de estudio considera los números en las diapositivas como un resumen de experiencia del conferenciantes en entornos de producción específicos. En particular, el \textbf{provider price}, la tasa de falla de las GPUs y la aceleración (\textbf{speedup}) deben interpretarse junto con la \textbf{workload}, el hardware, la versión del software, la línea base y la ventana de medición; sin estas condiciones, solo pueden servir como pistas de diseño y no como promesas universales extrapolables.
\end{warningbox}

\subsection{Resumen de la sección}

La dificultad de la inferencia proviene de las restricciones end a end: la aplicación define la corrección y la experiencia, y el sistema las convierte en problemas de recursos y agenda; el hardware y el software deciden juntos los límites factibles. Todas las técnicas posteriores obedecen el mismo orden: \textbf{definir primero la tarea y la medición, luego seleccionar la arquitectura y finalmente optimizar la implementación.}


